\chapter{Node-Timer und ABMs}
\label{ch:map/timers}

Eine Funktion regelmäßig auf bestimmten Nodes auszuführen, ist eine
häufige Aufgabe. Luanti bietet dafür zwei Wege: \emph{Active Block
Modifiers} (ABMs)\index{abm@ABM} und
Node-Timer.\index{node-timer@Node-Timer}

ABMs durchsuchen alle geladenen Mapblöcke nach Nodes, die bestimmte
Kriterien erfüllen. Sie eignen sich am besten für Nodes, die in der Welt
häufig vorkommen, etwa Gras. Sie kosten viel Rechenzeit, aber wenig
Arbeits- und Speicherplatz.

Für seltene Nodes oder solche, die ohnehin
Metadaten\index{metadaten@Metadaten} verwenden, etwa Öfen und Maschinen,
solltest du stattdessen Node-Timer verwenden. Node-Timer merken sich die
anstehenden Timer in jedem Mapblock\index{mapblock@Mapblock} und führen
sie aus, sobald sie ablaufen. Sie müssen also nicht alle geladenen Nodes
nach Treffern durchsuchen, brauchen dafür aber etwas mehr Arbeits- und
Speicherplatz, um die anstehenden Timer zu verwalten.

\section{Node-Timer}\label{node-timer}

Node-Timer sind direkt an einen einzelnen Node gebunden. Du steuerst sie
über ein \texttt{NodeTimerRef}-Objekt:

\begin{Shaded}
\begin{Highlighting}[]
\KeywordTok{local} \VariableTok{timer} \OperatorTok{=} \VariableTok{core}\OperatorTok{.}\NormalTok{get\_node\_timer}\OperatorTok{(}\VariableTok{pos}\OperatorTok{)}
\VariableTok{timer}\OperatorTok{:}\NormalTok{start}\OperatorTok{(}\DecValTok{10.5}\OperatorTok{)} \CommentTok{{-}{-} in Sekunden}
\end{Highlighting}
\end{Shaded}

Läuft ein Node-Timer ab, wird die Methode
\texttt{on\_\allowbreak{}timer}\index{callback@Callback!on_timer@\texttt{on\_\allowbreak{}timer}}
in der Definitionstabelle des Nodes aufgerufen. Sie bekommt nur einen
Parameter, die Position des Nodes:

\begin{Shaded}
\begin{Highlighting}[]
\VariableTok{core}\OperatorTok{.}\NormalTok{register\_node}\OperatorTok{(}\StringTok{"autodoors:door\_open"}\OperatorTok{,} \OperatorTok{\{}
    \VariableTok{on\_timer} \OperatorTok{=} \KeywordTok{function}\OperatorTok{(}\VariableTok{pos}\OperatorTok{)}
        \VariableTok{core}\OperatorTok{.}\NormalTok{set\_node}\OperatorTok{(}\VariableTok{pos}\OperatorTok{,} \OperatorTok{\{} \VariableTok{name} \OperatorTok{=} \StringTok{"autodoors:door"} \OperatorTok{\})}
        \ControlFlowTok{return} \KeywordTok{false}
    \KeywordTok{end}\OperatorTok{,}
\OperatorTok{\})}
\end{Highlighting}
\end{Shaded}

Gibt \texttt{on\_\allowbreak{}timer} \texttt{true} zurück, läuft der
Timer mit demselben Intervall erneut. Du kannst auch innerhalb von
\texttt{on\_\allowbreak{}timer}
\texttt{get\_\allowbreak{}node\_\allowbreak{}timer(pos)} verwenden; gib
dann aber \texttt{false} zurück, um Konflikte zu vermeiden.

Timer haben eine Einschränkung: Aus Optimierungsgründen gibt es je
Node-Typ nur eine Art von Timer, und je Node kann nur ein Timer laufen.

\section{Active Block Modifiers}\label{active-block-modifiers}

Alien-Gras ist für dieses Kapitel eine Grasart, die mit einer gewissen
Wahrscheinlichkeit in der Nähe von Wasser erscheint.

\begin{Shaded}
\begin{Highlighting}[]
\VariableTok{core}\OperatorTok{.}\NormalTok{register\_node}\OperatorTok{(}\StringTok{"aliens:grass"}\OperatorTok{,} \OperatorTok{\{}
    \VariableTok{description} \OperatorTok{=} \StringTok{"Alien{-}Gras"}\OperatorTok{,}
    \VariableTok{light\_source} \OperatorTok{=} \DecValTok{3}\OperatorTok{,} \CommentTok{{-}{-} Der Node leuchtet; min. 0, max. 14}
    \VariableTok{tiles} \OperatorTok{=} \OperatorTok{\{} \StringTok{"aliens\_grass.png"} \OperatorTok{\},}
    \VariableTok{groups} \OperatorTok{=} \OperatorTok{\{} \VariableTok{choppy} \OperatorTok{=} \DecValTok{1} \OperatorTok{\},}
    \VariableTok{on\_use} \OperatorTok{=} \VariableTok{core}\OperatorTok{.}\NormalTok{item\_eat}\OperatorTok{(}\DecValTok{20}\OperatorTok{),}
\OperatorTok{\})}

\VariableTok{core}\OperatorTok{.}\NormalTok{register\_abm}\OperatorTok{(\{}
    \VariableTok{nodenames} \OperatorTok{=} \OperatorTok{\{} \StringTok{"default:dirt\_with\_grass"} \OperatorTok{\},}
    \VariableTok{neighbors} \OperatorTok{=} \OperatorTok{\{}
        \StringTok{"default:water\_source"}\OperatorTok{,}
        \StringTok{"default:water\_flowing"}\OperatorTok{,}
    \OperatorTok{\},}
    \VariableTok{interval} \OperatorTok{=} \DecValTok{10.0}\OperatorTok{,} \CommentTok{{-}{-} alle 10 Sekunden ausführen}
    \VariableTok{chance} \OperatorTok{=} \DecValTok{50}\OperatorTok{,} \CommentTok{{-}{-} jeder Node wird mit Chance 1:50 gewählt}
    \VariableTok{action} \OperatorTok{=} \KeywordTok{function}\OperatorTok{(}
        \VariableTok{pos}\OperatorTok{,}
        \VariableTok{node}\OperatorTok{,}
        \VariableTok{active\_object\_count}\OperatorTok{,}
        \VariableTok{active\_object\_count\_wider}
    \OperatorTok{)}
        \KeywordTok{local} \VariableTok{above} \OperatorTok{=}
            \OperatorTok{\{} \VariableTok{x} \OperatorTok{=} \VariableTok{pos}\OperatorTok{.}\VariableTok{x}\OperatorTok{,} \VariableTok{y} \OperatorTok{=} \VariableTok{pos}\OperatorTok{.}\VariableTok{y} \OperatorTok{+} \DecValTok{1}\OperatorTok{,} \VariableTok{z} \OperatorTok{=} \VariableTok{pos}\OperatorTok{.}\VariableTok{z} \OperatorTok{\}}
        \VariableTok{core}\OperatorTok{.}\NormalTok{set\_node}\OperatorTok{(}\VariableTok{above}\OperatorTok{,} \OperatorTok{\{} \VariableTok{name} \OperatorTok{=} \StringTok{"aliens:grass"} \OperatorTok{\})}
    \KeywordTok{end}\OperatorTok{,}
\OperatorTok{\})}
\end{Highlighting}
\end{Shaded}

Dieser ABM läuft alle zehn Sekunden, und für jeden passenden Node wird
er mit einer Wahrscheinlichkeit von 1 zu 50 ausgeführt. Läuft der ABM
auf einem Node, wird darüber ein Alien-Gras-Node gesetzt. Achtung: Dabei
wird jeder Node gelöscht, der sich vorher an dieser Stelle befand. Um
das zu verhindern, solltest du mit
\texttt{core.\allowbreak{}get\_\allowbreak{}node}\index{core@\texttt{core}!get_node@\texttt{core.\allowbreak{}get\_\allowbreak{}node}}
prüfen, ob dort Platz für das Gras ist.

Die Angabe von Nachbarn (\texttt{neighbors}) ist optional. Gibst du
mehrere Nachbarn an, genügt es, wenn einer davon vorhanden ist.

Auch \texttt{chance} ist optional. Ohne Angabe läuft der ABM immer, wenn
die übrigen Bedingungen erfüllt sind.

\section{Jetzt bist du dran}\label{jetzt-bist-du-dran}

\begin{itemize}
\tightlist
\item
  Midas-Berührung: Lass Wasser alle 5 Sekunden mit einer
  Wahrscheinlichkeit von 1 zu 100 zu Goldblöcken werden.
\item
  Verfall: Lass Holz zu Erde werden, wenn Wasser angrenzt.
\item
  Lauffeuer: Lass jeden Luft-Node in Flammen aufgehen. (Tipp:
  \texttt{\textquotedbl{}air\textquotedbl{}} und
  \texttt{\textquotedbl{}fire:basic\_\allowbreak{}flame\textquotedbl{}}.)
  Warnung: Rechne damit, dass das Spiel abstürzt.
\end{itemize}
